\begin{proof}
\textbf{Statement.}  
Let  
\[
a_{1}\le a_{2}\le\cdots\le a_{n},\qquad 
b_{1}\le b_{2}\le\cdots\le b_{n}
\]
be real numbers.  We must show
\begin{align}
n\sum_{i=1}^{n}a_{i}b_{i}\;\ge\;
\Bigl(\sum_{i=1}^{n}a_{i}\Bigr)\Bigl(\sum_{i=1}^{n}b_{i}\Bigr).\tag{1}
\end{align}

\medskip\noindent\textbf{1. Trivial and degenerate cases.}
\begin{itemize}
\item If \(n=1\), (1) reduces to \(a_{1}b_{1}\ge a_{1}b_{1}\), true with equality.
\item If all \(a_{i}\) are equal, say \(a_{i}=a\), then  
\(\sum a_{i}b_{i}=a\sum b_{i}\) and both sides of (1) equal \(na\sum b_{i}\).
\item The case “all \(b_{i}\) equal’’ is symmetric.
\end{itemize}
Hence we may assume \(n\ge2\) and that neither sequence is constant.

\medskip\noindent\textbf{2. Rearrangement Inequality.}  
For increasing sequences \(x_{1}\le\cdots\le x_{n}\) and
\(y_{1}\le\cdots\le y_{n}\) and any permutation \(\sigma\in S_{n}\),
\[
\sum_{i=1}^{n}x_{i}y_{\sigma(i)}\le
\sum_{i=1}^{n}x_{i}y_{i}\le
\sum_{i=1}^{n}x_{i}y_{\,n+1-i}.
\]
Thus the sum with the identity permutation is maximal.  
Applying this with \(x_{i}=a_{i}\) and \(y_{i}=b_{i}\) yields
\begin{align}
\sum_{i=1}^{n}a_{i}b_{i}
\text{ is the largest of the } n!
\text{ numbers }
\sum_{i=1}^{n}a_{i}b_{\sigma(i)}\;( \sigma\in S_{n}).\tag{2}
\end{align}

\medskip\noindent\textbf{3. Average over all permutations.}  
Consider
\begin{align}
S:=\sum_{\sigma\in S_{n}}\;\sum_{i=1}^{n}a_{i}\,b_{\sigma(i)}.\tag{3}
\end{align}
Fix an index \(i\).  In the collection of all permutations each \(b_{j}\) appears
in position \(i\) exactly \((n-1)!\) times, hence
\[
\sum_{\sigma\in S_{n}}b_{\sigma(i)}=(n-1)!\sum_{j=1}^{n}b_{j}.
\]
Multiplying by \(a_{i}\) and summing over \(i\) gives
\begin{align}
S=(n-1)!\Bigl(\sum_{i=1}^{n}a_{i}\Bigr)\Bigl(\sum_{j=1}^{n}b_{j}\Bigr).\tag{4}
\end{align}
Dividing by the number of permutations \(n!\) we obtain the average value:
\begin{align}
\frac{S}{n!}= \frac{1}{n}\Bigl(\sum_{i=1}^{n}a_{i}\Bigr)\Bigl(\sum_{j=1}^{n}b_{j}\Bigr).\tag{5}
\end{align}

\medskip\noindent\textbf{4. Comparison with the maximal sum.}  
From (2) the ordered sum \(\sum a_{i}b_{i}\) is not smaller than any other
\(\sum a_{i}b_{\sigma(i)}\); therefore it is not smaller than their average.  Using (5),
\[
\sum_{i=1}^{n}a_{i}b_{i}\;\ge\;
\frac{S}{n!}
=
\frac{1}{n}\Bigl(\sum_{i=1}^{n}a_{i}\Bigr)\Bigl(\sum_{j=1}^{n}b_{j}\Bigr).
\]
Multiplying by \(n\) yields exactly inequality (1).

\medskip\noindent\textbf{5. Equality condition.}  
Equality in the Rearrangement Inequality (2) occurs iff the two sequences are
not strictly ordered with respect to each other, i.e. at least one of the
sequences is constant.  Consequently equality in Chebyshev’s inequality (1)
holds precisely when all \(a_{i}\) are equal or all \(b_{i}\) are equal – the
situations already treated in Step 1. No other case can give equality, because
then the ordered sum would be strictly larger than the average in Step 4.

\medskip\noindent\textbf{6. Conclusion.}  
Using the Rearrangement Inequality together with a simple counting argument we
have proved that for any two similarly ordered real sequences
\[
a_{1}\le a_{2}\le\cdots\le a_{n},\qquad
b_{1}\le b_{2}\le\cdots\le b_{n},
\]
the Chebyshev sum inequality
\[
n\sum_{i=1}^{n}a_{i}b_{i}\ge
\Bigl(\sum_{i=1}^{n}a_{i}\Bigr)\Bigl(\sum_{i=1}^{n}b_{i}\Bigr)
\]
holds, with equality exactly when one of the sequences is constant. \(\square\)
\end{proof}